\documentclass[../main]{subfiles}
\begin{document}
\chapter{Nuevos Materiales}
\img{images/13/grafeno.jpg}{Nuevos Materiales.}

\info{Nuevos Materiales}{Los nuevos materiales son el resultado del desarrollo científico y tecnológico actual. Su creación se basa en el conocimiento de las moléculas y los átomos para obtener propiedades específicas, proceso impulsado fundamentalmente por la Nanotecnología.}

\section{Guía de Propiedades de los Materiales}
Para seleccionar un material para un componente industrial, es necesario analizar tres dimensiones fundamentales:

\subsection{Propiedades Mecánicas}
Definen cómo reacciona el material ante fuerzas físicas:
\begin{itemize}
  \item \textbf{Resistencia a la Tracción}: Es la máxima tensión que el material puede soportar antes de romperse.
  \item \textbf{Dureza}: Resistencia a la indentación superficial o al rayado.
  \item \textbf{Elasticidad (Módulo de Young)}: Capacidad de recuperar la forma original tras una deformación elástica.
\end{itemize}

\subsection{Propiedades Térmicas}
Definen la reacción del material ante el calor:
\begin{itemize}
  \item \textbf{Conductividad Térmica}: Capacidad de transferir calor (alta en disipadores, baja en aislantes).
  \item \textbf{Coeficiente de Expansión Térmica}: Variación dimensional al calentarse (crítico para ajustes de precisión).
  \item \textbf{Punto de Fusión}: Temperatura a la cual el material pasa al estado líquido.
\end{itemize}

\subsection{Propiedades Químicas}
Definen la interacción con el entorno:
\begin{itemize}
  \item \textbf{Resistencia a la Corrosión}: Capacidad de resistir la oxidación (ej. Acero Inoxidable vs. Acero al Carbono).
  \item \textbf{Inercia Química}: Ausencia de reactividad frente a ácidos o bases.
\end{itemize}

\section{Tutorial: Matriz de Comparación de Materiales}
Cuando el proyecto requiere una propiedad específica (por ejemplo, "debe ser ligero pero resistente"), se debe seguir este proceso sistemático:

\begin{enumerate}
  \item \textbf{Listar Requerimientos}: Definir límites cuantificables (ej. Peso $<$ 2kg, Resistencia $>$ 500MPa, Costo $<$ \$100).
  \item \textbf{Identificar Candidatos}: Seleccionar materiales potenciales (ej. Aluminio 6061, Fibra de Carbono, Titanio Grado 5).
  \item \textbf{Crear Matriz de Comparación}: Calificar cada material frente a los requerimientos establecidos.
  \item \textbf{Análisis de Compromisos (Trade-off)}: Si ningún material cumple todos los criterios, decidir qué requerimiento puede flexibilizarse (ej. aceptar un mayor costo para obtener el peso requerido).
\end{enumerate}

\section{Guía de Planificación de I+D para Ensayos de Materiales}
Al testear un nuevo material (como un compuesto impreso en 3D), se recomienda la siguiente metodología de trabajo:

\subsection{Fase 1: Hitos (Milestones)}
Definir el objetivo claro del ensayo. \textit{Ejemplo: "Sustituir el soporte de acero por un polímero reforzado para reducir el peso en un 30\%"}.

\subsection{Fase 2: Caracterización}
Realizar pruebas controladas sobre muestras del material:
\begin{itemize}
  \item Ensayos de tracción y dureza.
  \item Ciclos térmicos para evaluar la estabilidad.
  \item Pruebas de fatiga.
\end{itemize}

\subsection{Fase 3: Viabilidad}
Analizar la relación costo de producción vs. ganancia de rendimiento. Verificar la disponibilidad del material en el mercado local.

\subsection{Fase 4: Sostenibilidad}
Evaluar el impacto ambiental: ¿Es reciclable? ¿El proceso de producción es tóxico?

\end{document}